% page 20 of the facsimile (image p-019 of the scan)
% block 157.5 x 220.4 mm ; left margin 8.0 mm
\pgc{-12.700mm}{0.000mm}{
\l{\cel{5}12}
\l{}
\l{\sou{-ajo}. Sufixo qua indikas : I. Kozo facita ek ula materio,}
\l{\cel{3}o qua posedas ta o ca karaktero (Ex.: lanajo, kotonajo,}
\l{\cel{3}bonajo, bitrajo). - II. (\sou{per extenso di la senco}). Ago,}
\l{\cel{3}parolo, procedo di... (amikalajo, pueralajo). - III.}
\l{\cel{3}(\sou{kun verbo transitiva}) Indikas la objekto pasiva di la}
\l{\cel{3}ago (cakaze, lu equivalas \textquotedbl{}-ata\textquotedbl{}) (Ex.: manjajo :}
\l{\cel{3}karno, pano, legumi) ; vidajo : kozo vidata. - IV. (\sou{Kun}}
\l{\cel{3}\sou{verbo mixita})\cel{1}\sou{-ajo}\cel{1}=\cel{1}\sou{a(ta)jo}\cel{1}(Ex.: chanjajo : kozo chan-}
\l{\cel{3}jata.)\cel{2}V. (\sou{Kun verbo netransitiva}). Indikas la subjekto,}
\l{\cel{3}ed equivalas\cel{2}\sou{-(ant)ajo}\cel{1}(Ex. : rezultajo : to quo rezul-}
\l{\cel{3}tas ; eventajo : to quo eventas).}
\l{}
\l{\cel{1}\sou{ajio}. (\sou{banko-komerco})Profito de la difero inter la valoro}
\l{\cel{3}nominala e la valoro reala di la moneto-peci - inter la}
\l{\cel{3}valoro nuna e la preco lor la pagoteso di la valor-paperi,}
\l{\cel{3}di la vari di qui la kurso varias. - DEFIRS.}
\l{}
\l{\sou{ajila}.\cel{2}Di qua la movi esas flexebla, rapida. - EFIS.}
\l{}
\l{\sou{ajiotar}. (\sou{netrans., pri)}\cel{2}Spekular pri la valor-paperi, pri}
\l{\cel{3}la vari, di qui la kursi esas variiva. - DEFIRS.}
\l{}
\l{\sou{ajornar}.\cel{2}(\sou{trans., til}). Rezolvar ke (ulo) agesos o facesos}
\l{\cel{3}erste en ta o ca dio, vice nun. - EFI.}
\l{}
\l{\sou{ajuro}. I. (\sou{arkitekt.}) Singla ek la mikra aperturi di ornament-}
\l{\cel{3}ajo. - II. (\sou{stofo}) Quaza truo en broduro o dentelo.- DEF.}
\l{}
\l{\sou{ajustar}. (\sou{trans.}) Intermuntar ed inter-harmoniigar la parti}
\l{\cel{3}di mashino, di mekanismo, por atingar la skopo determinita.}
\l{\cel{3}- EF.}
\l{}
\l{\sou{ajuto}. Tubo adaptita a la orifico di fonteno, di recipiento,}
\l{\cel{3}e c. por modifikar (pri lua volumo o lua formo) la fluo di}
\l{\cel{3}la liquido. - EF.}
\l{}
\l{\sou{akacio}. (\sou{bot.}) Genero de planti mimozea, di qui ula speci}
\l{\cel{3}donas kachuo, glutino-gumo, e c.\cel{2}L.\cel{1}\sou{acacia vera}.- DEFIRS.}
\l{}
\l{\sou{akademio}. Societo literaturala, ciencala, artala, di qua la}
\l{\cel{3}membri esas, generale, elektita e fixa-nombra, e qua ne}
\l{\cel{3}docas. - DEFIRS.}
\l{}
\l{\sou{akanto}.\cel{1}\sou{(bot}.)Planto bikotiledona, di qua la folii esas dent-}
\l{\cel{3}oza. - L.\cel{1}\sou{acanthus}. - DEFIRS.}
\l{}
\l{.\sou{akantino.}\cel{1}(\sou{biol}.)Substanco organika, albuminoza, solvebla}
\l{\cel{3}en sal-aquo - ye qua konsistas la skeleto di ula radio-}
\l{\cel{3}larii.}
}
